% Amostra do estilo: um capitulo de teste com todos os elementos do livro.
% Serve para conferir mudancas no estilo/livro.sty antes de aplica-las:
%   python compilar.py amostra             (versao rascunho, com notas)
%   python compilar.py amostra --modo final
% O texto abaixo e so enchimento para testar a diagramacao, nao e conteudo do livro.
\documentclass[11pt,twoside,openright,showtrims]{memoir}
\usepackage{estilo/livro}
% Dados do livro, num lugar so. A capa interna, a pagina de creditos, os
% cabecalhos e as propriedades do PDF leem daqui.

\newcommand{\titulolivro}{Do Cubo ao Copiloto}% titulo de trabalho, ver notas/02-titulos.md
\newcommand{\subtitulolivro}{Análise e engenharia de dados antes e depois da inteligência artificial}
\newcommand{\autorlivro}{Felipe Ramires Terrazas}
\newcommand{\edicaolivro}{1ª edição}
\newcommand{\anolivro}{2027}
\newcommand{\cidadelivro}{\pendente{cidade}}
\newcommand{\editoralivro}{Edição do autor}
\newcommand{\isbnlivro}{\pendente{ISBN}}

\addbibresource{referencias.bib}

\begin{document}
\mainmatter
\part{Parte de teste}
\label{parte:teste}

\chapter{Um capítulo de amostra}
\label{cap:amostra}

\epigraph{Uma epígrafe curta na abertura do capítulo, em itálico e alinhada à direita.}{Fonte da epígrafe}

\rascunho{Nota do autor: aparece só no PDF de rascunho. Serve para planejar o capítulo
sem misturar o plano com o texto.}

Este parágrafo de abertura testa a fonte do texto corrido, o espaçamento entre linhas
e a largura da mancha. Uma linha confortável de livro tem entre 60 e 75 caracteres,
e o formato 16 por 23 centímetros, com estas margens, fica nessa faixa. Uma citação no
texto aparece assim \parencite{kimball2013}, e uma nota de rodapé, assim.\footnote{A nota
de rodapé usa o corpo menor e fica separada do texto por um fio curto.}

O segundo parágrafo começa recuado. Termos de código aparecem no meio da frase com
\codigo{SELECT}, \codigo{GROUP BY} ou o nome de uma tabela como \codigo{fato\_vendas}.
Uma remissão para outra seção sem número fica assim: \veja{sec:amostra-codigo}.
Um ponto ainda sem resposta fica marcado \pendente{exemplo de pendência}.

\section{Uma seção com caixas}
\label{sec:amostra-caixas}

O texto segue normalmente antes da caixa de relato, que é o lugar das histórias de
projetos reais.

\begin{relato}
Aqui entra um relato de um projeto real, sem nomes de empresa nem de cliente. A caixa
tem fundo cinza claro (para gastar pouca tinta) e um fio à esquerda. Ela pode
quebrar entre páginas quando o relato for longo.
\end{relato}

\begin{destaque}
Uma ideia-chave em uma ou duas frases, para o leitor que folheia o livro.
\end{destaque}

\begin{dica}
Uma dica prática, curta, que o leitor pode aplicar no trabalho da semana seguinte.
\end{dica}

\antesdepois{O mesmo trabalho como era feito antes de a IA entrar no fluxo.}
{O mesmo trabalho hoje, com o que mudou e o que continuou igual.}

\subsection{Uma subseção}

Texto da subseção, para conferir a hierarquia dos títulos.

\section{Uma seção com código e tabela}
\label{sec:amostra-codigo}

\begin{sql}
-- Receita por mês e por região (comentário com acentuação)
SELECT
    d.ano_mes,
    r.regiao,
    SUM(f.valor) AS receita
FROM fato_vendas AS f
JOIN dim_data   AS d ON d.id_data   = f.id_data
JOIN dim_regiao AS r ON r.id_regiao = f.id_regiao
GROUP BY d.ano_mes, r.regiao
ORDER BY d.ano_mes;
\end{sql}

\begin{python}
import pandas as pd

vendas = pd.read_parquet("vendas.parquet")
por_mes = vendas.groupby("ano_mes")["valor"].sum()  # soma por mês
\end{python}

\begin{table}[htbp]
  \centering
  \caption{Uma tabela no padrão do livro.}
  \label{tab:amostra}
  \begin{tabular}{@{}lrr@{}}
    \toprule
    Região & Receita & Variação \\
    \midrule
    Norte  & 1.250,00 & 4,2\% \\
    Sul    & 3.980,50 & $-$1,7\% \\
    \bottomrule
  \end{tabular}
\end{table}

A Tabela~\ref{tab:amostra} mostra o padrão: fios do \codigo{booktabs}, sem linhas
verticais, números alinhados à direita.

\printbibliography[heading=bibintoc,title={Referências}]
\end{document}
